\itemhead{复习题三 · 第13题（教材p.91）}{XJ-B1-C03-RV-P091-Y0785}
\label{XJ-B1-C03-RV-P091-Y0785}
{\footnotesize 湘教版教材p.91 / PDF第98页 / 复习题三 / 第13题}
\field{原题}{根据如图所示的图象，简要说明近 150 年来人类消耗的能源结构变化情况，并对未来 100 年能源结构的变化趋势作出预测。
\bankfig{C03-RV091}{0.7}}
\field{参考答案}{按所给历史示意图：木材比例总体下降，煤炭先升后降，石油先升后趋降，天然气和核能从出现后逐渐增加。未来趋势只能作基于此图的推测，可讨论核能、水合物和太阳能的份额增加；这不是现实能源统计或确定预测，不能用示意图编造实际数据。}
{\small \field{关联知识}{F04 函数三种表示法、图象与读数（教材帮p.85、86、93、94 / PDF94、95、102、103）}
\field{方法与易错点}{把关键点与定义域一起标到图上，逐区间读交点和上下关系。 易错：将示意图预测当事实；把两曲线交点收进严格不等式；折线漏端点。}
\field{答案依据}{答案册相关题组：PDF第24、25、26、27页（册内页23、24、25、26）；按当前题库原题逐项重算，不按旧版题号机械对套。}}
